\documentclass[11pt,a4paper]{article}
\usepackage[margin=2cm]{geometry}
\usepackage{booktabs,array,longtable}
\usepackage[table]{xcolor}
\usepackage{titlesec}
\usepackage{enumitem}
\usepackage[colorlinks,allcolors=blue!55!black]{hyperref}
\usepackage{parskip}
\definecolor{hd}{RGB}{18,48,88}
\definecolor{bad}{RGB}{150,20,20}
\definecolor{good}{RGB}{20,110,45}
\titleformat{\section}{\large\bfseries\color{hd}}{\thesection}{0.6em}{}
\titleformat{\subsection}{\normalsize\bfseries\color{hd}}{\thesubsection}{0.6em}{}
\newcommand{\B}[1]{\textcolor{bad}{\textbf{#1}}}
\newcommand{\G}[1]{\textcolor{good}{\textbf{#1}}}
\setlist{itemsep=2pt,topsep=3pt,parsep=0pt}
\renewcommand{\arraystretch}{1.2}
\title{\vspace{-1.4cm}\textbf{COMPASS Data Collection Plan}\\[1.5mm]
\large Three dense reference floors across two buildings\\
\normalsize Seven phones, cart-mounted survey, existing logging application}
\date{8 September 2026}\author{}
\begin{document}\maketitle\vspace{-8mm}\hrule\vspace{3mm}

\section{Objective}

Produce the \textbf{largest publicly released measured indoor radio map with surveyed
transmitter positions}, with multi-device measurements and inertial streams.

\subsection{Novelty claim, audited 8 September 2026}

An earlier draft claimed the \emph{first} dense ground-truth indoor radio map. \B{That
claim is false} and must not be used. A literature audit found four measured indoor
counterexamples, and one of them already claims a ``first'' of its own.

\begin{center}\small
\begin{tabular}{@{}p{4.3cm}p{5.0cm}p{6.4cm}@{}}
\toprule
Counterexample & What it has & Where it loses to us \\ \midrule
Milosheski et al., \emph{Data in Brief} 2026 (arXiv 2511.00494) & Measured, indoor, WiFi,
\textbf{20 surveyed AP positions published}, 53 points at 2\,m & 53 points over
$\sim$212\,m$^2$, one device, one building, no inertial data \\
OpenCSI (arXiv 2104.07963) & Measured LTE CSI, 1\,cm spacing, surveyed ceiling
transmitter. Already claims ``first publicly available radio map containing CSI'' &
17.5\,m$^2$ total, one transmitter, sampled along 8 lines rather than areally, CSI not
calibrated RSS \\
DICHASUS, KU~Leuven MaMIMO & 178k--427k points, millimetre position accuracy &
The \textbf{surveyed nodes are receivers}, transmitter is mobile; CSI/SNR, not pathloss;
KU~Leuven covers $1.25\times1.25$\,m \\
NeRF$^2$ BLE (MobiCom 2023) & $\sim$6000 points over $\sim$1394\,m$^2$, 50 surveyed
gateways & Gateways are receivers; uncalibrated commodity BLE RSSI \\
Romero et al.\ (arXiv 2310.11036) & Claims ``first public RME datasets with real data'' &
\textbf{Outdoor} \\
\bottomrule
\end{tabular}
\end{center}

\textbf{What does survive, and it is the useful half.} The indoor radio-map benchmark
canon is entirely simulated. The ICASSP~2025 and MLSP~2025 Indoor Radio Map Dataset is
\G{ray-traced with Ranplan Wireless} at 0.25\,m, confirmed from the challenge site, the
overview paper and the IEEE DataPort record. A March 2026 tutorial surveying 16 public
radio-map datasets finds only two involve measurement at all, and \emph{every} 2-D indoor
pathloss dataset is simulated.

\textbf{Claim to use.} ``Indoor radio-map benchmarks are ray-traced. Public measured
indoor releases are either sparse (53 points at 2\,m), confined to a few square metres,
or publish receiver rather than transmitter geometry. We release $\sim$820 measured
points at 1.5\,m across three floors in two buildings, with surveyed transmitters, seven
devices and inertial streams.'' Every clause there is individually citable.

\textbf{Must cite.} Milosheski et al., OpenCSI, NeRF$^2$, DICHASUS and Romero et al. A
reviewer who finds any of these uncited will assume we did not look.

The existing UniCellular corpus is already published, so it cannot itself be our
contribution. What we collect now can be. The scientific target is not more data but a
different \emph{kind} of data: a field dense enough to hold out contiguous regions
rather than scattered points.

\subsection{The defect this fixes}

Reference points currently sit \textbf{3.8 to 4.1\,m apart}. At that spacing the only
possible evaluation is to hide one point and predict it from its neighbours, which asks
solely for interpolation between nearby samples of one smooth process. That is the exact
regime where a tuned interpolator is near-optimal by construction, and it is why
classical methods beat every learned method on our real data today.

The protocol is provably degenerate: the 0\,m and 1\,m exclusion buffers currently
produce \B{byte-identical} results, because no two reference points are ever within a
metre of each other.

With a dense field we can mask a contiguous block, a whole room, or everything beyond a
given wall. An interpolator then has nothing to interpolate from; a geometry-conditioned
model can still reason. \textbf{That is the comparison our method exists for, and the one
the current data physically cannot express.}

\section{Scope decisions}

\subsection{In: three floors, two buildings}

\begin{center}\small
\begin{tabular}{@{}llrrl@{}}
\toprule
Floor & Building type & Extent (m) & Current RPs & Role \\ \midrule
\rowcolor{green!8}CMU-Q f2 & Walled academic & $39.7\times30.2$ & 84 & Primary dense reference \\
CMU-Q f1 & Walled academic & $41.1\times29.2$ & 84 & Cross-floor axis, same building \\
EC Parking f1 & Open concrete & $48.6\times11.8$ & 34 & \textbf{Negative control} \\
\bottomrule
\end{tabular}
\end{center}

All three already have floor plans on disk. CMU-Q f2 is primary because it has the most
transmitters (35), the widest device spread in the corpus (13.5\,dB) and six of seven
phones sharing cells.

\textbf{Why the control site matters more than a third building.} Wall-aware
interpolation currently gives $+3.86\%$ at CMU-Q and \G{exactly $0.00\%$} at EC Parking,
where cross-validation independently selects $\lambda=0$ at both folds. A geometric prior
that helps where walls exist and does nothing where they do not is far stronger evidence
than three sites that all agree.

\subsection{Out: both Ezdan towers}

Not because floor plans are missing, though they are. Every Ezdan floor reports an extent
of roughly $25.6\times\B{0.0}$\,m. \textbf{The y-extent is zero on all 55 floors}, meaning
every reference point on a floor shares one y coordinate and the points lie on a single
line. That is 1-D data, and 2-D radio-map reconstruction on a line is not defined. With
2--5 points per floor and zero evaluable cells, Ezdan is excluded on geometry.

\section{Equipment}

\begin{center}\small
\begin{tabular}{@{}p{3.6cm}p{12.4cm}@{}}
\toprule
Item & Specification and rationale \\ \midrule
7 phones & Mounted in \textbf{fixed relative geometry} on the cart. Fixed geometry is what
makes per-point device comparison controlled rather than confounded by position. Record
the physical layout once and never change it mid-campaign. \\
Survey cart & Carries all seven phones at a \textbf{fixed height of 1.2--1.5\,m},
approximating hand-held height. Constant height matters: antenna height changes received
power, and a varying height would alias into the spatial signal we are trying to measure. \\
Logging application & Existing. See the capability checklist in Section~6 for the four
features the protocol depends on. \\
Laser rangefinder & For grid layout and for verifying plan scale. \\
Chalk or removable floor markers & Roughly 320 marks per floor. \\
\bottomrule
\end{tabular}
\end{center}

\subsection{Cart for the survey, hands for the trajectories}

The cart is right for the dense grid because it gives precision, repeatability and
constant geometry. It is \emph{wrong} for the trajectory sessions, because a cart is not
how anyone carries a phone, and body attenuation and carry mode are part of the science.

\textbf{Use the cart for Steps 2 and 3. Use human carriage for Step 4.} Do not substitute
one for the other.

\section{Protocol}

\subsection{Step 0 --- Verify plan scale (half a day, all floors)}

Scale is solved for CMU-Q and EC Parking, so this is verification rather than derivation,
but it must still be recorded because every geometric result depends on it.

\begin{enumerate}
\item Measure two long known distances per floor with the rangefinder: one full corridor
run, and one perpendicular span.
\item Locate the same two features on the plan PNG in pixels.
\item Derive metres per pixel \textbf{per floor and per axis}.
\item If the x and y values disagree, the plan is anisotropic and must be rectified
before use.
\end{enumerate}

\textbf{Record the result in the released metadata.} The current code applies a single
global constant of 0.032\,m/px to every site, and all four plans are 1650$\times$1169
pixels despite covering different real areas, so a per-floor number is a genuine
correction rather than bookkeeping.

\subsection{Step 1 --- Lay the grid (half a day per floor)}

\begin{enumerate}
\item Establish the frame from two perpendicular walls.
\item Mark a regular \textbf{1.5\,m} grid across all walkable space.
\item Target under \textbf{0.3\,m} positional error, one fifth of the pitch.
\item Number every mark and record its coordinate in the frame.
\end{enumerate}

\begin{center}\small
\begin{tabular}{@{}lrrrr@{}}
\toprule
Floor & Gross area & Walkable (est.) & Grid points at 1.5\,m & Versus today \\ \midrule
CMU-Q f2 & 1199\,m$^2$ & $\sim$700\,m$^2$ & $\sim$310 & 3.7$\times$ \\
CMU-Q f1 & 1200\,m$^2$ & $\sim$700\,m$^2$ & $\sim$310 & 3.7$\times$ \\
EC Parking f1 & 573\,m$^2$ & $\sim$450\,m$^2$ & $\sim$200 & 5.9$\times$ \\
\midrule
\textbf{Total} & & & \textbf{$\sim$820} & \\
\bottomrule
\end{tabular}
\end{center}

Why 1.5\,m: at 2.0\,m the floor yields only 175 points, which is the bare minimum for
block hold-out. At 1.0\,m it yields 700 and costs three to four days. 1.5\,m is the knee.

\subsection{Step 2 --- Cart survey of the grid (1.5 days per floor)}

At every mark:

\begin{enumerate}
\item Position the cart on the mark.
\item \textbf{Dwell 45 seconds.}
\item Log every visible cell on all seven phones, and \B{log the non-detections too}.
A phone failing to hear a cell is a measurement, not a gap. The detection-heterogeneity
result depends entirely on it: cross-device heard-set overlap is 0.299, and binary
heard-set matching (8.41\,m) \emph{beats} signal matching (9.25\,m). That finding is only
computable if absences are recorded.
\item Record the cart \textbf{heading}.
\item At every \textbf{tenth mark}, repeat the measurement at a second heading, so
orientation effects are quantified rather than assumed away.
\end{enumerate}

Timing: 45\,s dwell plus roughly 20\,s to move and settle the cart is 65\,s per point.
310 points is about 5.6 hours of collection.

\subsection{Step 3 --- Transmitter survey (half a day per floor)}

This is new work at every site. No public radio-map dataset has it, and \textbf{without it
the occlusion feature cannot be computed on real data at all}, because the feature works
by casting a ray from a source.

\begin{enumerate}
\item Physically locate every reachable antenna, femtocell, booster and repeater.
\item Record each position in the Step 0 frame, plus \textbf{mounting height}.
\item Record \textbf{which cell identifiers each physical unit broadcasts}.
\item Photograph each unit with its label.
\end{enumerate}

Twenty units per floor is sufficient. Currently only \B{17.9\%} of cells yield a
locatable source by signal fitting, with no external reference to validate against. This
step converts ``boosters share cell identifiers'' from an inference into a published fact.

\subsection{Step 4 --- Human-carried trajectories (1 day per floor)}

Existing mobile walks carry \texttt{rpNumber} $=-1$. They have \B{no position labels at
all}, so every order-awareness and active-sensing result is computed on sequences whose
location is unknown.

\begin{enumerate}
\item Walk routes that start, end and pass through marked grid points, pressing the app's
\textbf{marker button} at each, so position between fixes is interpolated from known
anchors rather than guessed.
\item Four route types per floor:
\begin{itemize}
\item corridor loop
\item a route crossing at least four walls
\item a route into a deep non-line-of-sight pocket
\item an open-space control
\end{itemize}
\item Three phones carried together per walk, in three carry modes: \textbf{hand, pocket,
bag}, logged per session.
\item Three times of day, three separate days.
\item Vary walking speed deliberately across repeats: slow, normal, brisk. Speed aliasing
is an unmeasured confound in the order argument.
\end{enumerate}

\subsection{Step 5 --- Device calibration capture (2 hours per campaign)}

All seven phones on the cart, same orientation, two-minute dwell, at \textbf{five points
spanning strong to weak signal}. Repeat at the start and the end of the campaign to
capture drift.

This yields a ground-truth offset \textbf{as a function of signal level}. It matters
because the current model assumes one additive scalar per device, while offsets are
estimable at 0.78--0.89 correlation yet calibration buys only 0.18--2.04\% in
reconstruction, which suggests the additive assumption is wrong.

\subsection{Step 6 --- Clock discipline (every session)}

\begin{enumerate}
\item Synchronise all phones to one NTP source at session start.
\item Log both wall-clock and monotonic time per sample.
\item Begin every session with a shared physical event, all phones held together for ten
seconds, so residual offset is recoverable in post-processing.
\end{enumerate}

\section{Application capability checklist}

The app exists, so this is a confirmation list rather than a build. Four features the
protocol depends on:

\begin{center}\small
\begin{tabular}{@{}p{4.6cm}p{11.4cm}@{}}
\toprule
Capability & Why the protocol needs it \\ \midrule
\textbf{Non-detection logging} & Must record cells that were scanned for and \emph{not}
heard, not only the ones received. Without this the detection-heterogeneity result cannot
be computed. If the app only stores received cells, this is the one change worth making
before collection starts. \\
\textbf{Marker / waypoint button} & Stamps a known ground-truth position mid-walk so
trajectory position is interpolable. \\
\textbf{Inertial streams} & Accelerometer and gyroscope at $\ge$50\,Hz, magnetometer at
$\ge$25\,Hz, barometer at $\ge$1\,Hz, orientation quaternion, step counter. No radio-map
dataset currently carries these, so this is a differentiator as well as a requirement. \\
\textbf{Dual clocks} & Monotonic alongside wall-clock, so cross-device drift is
recoverable after the fact. \\
\bottomrule
\end{tabular}
\end{center}

\section{What ships}

The contribution is a packaged artifact, not raw logs.

\begin{enumerate}
\item Dense reference fields for three floors, seven devices per point, \textbf{including
non-detections}
\item Surveyed transmitter positions, heights, and the physical-unit to cell-identifier
mapping
\item Located trajectories with inertial streams and carry-mode labels
\item Floor plans as occupancy rasters in the calibrated frame, with \textbf{per-floor,
per-axis} metres-per-pixel recorded
\item Per-device calibration curves as a function of signal level
\item \textbf{Frozen evaluation splits}, including the block hold-out region definitions
\item A datasheet documenting collection, devices, coverage and limitations
\item Baseline results for the 18 methods already implemented
\end{enumerate}

Item 6 is the one most often omitted and the one that makes a benchmark comparable rather
than merely available. The repository currently has no split manifest at all.

\section{Timeline}

\begin{center}\small
\begin{tabular}{@{}llr@{}}
\toprule
Phase & Content & Days \\ \midrule
App confirmation & Verify the four capabilities, pilot one corridor end to end & 1 \\
Scale verification & Step 0, all three floors & 0.5 \\
CMU-Q f2 & Steps 1--4 & 3 \\
CMU-Q f1 & Steps 1--4 & 3 \\
EC Parking f1 & Steps 1--4, smaller floor & 2.5 \\
Calibration & Step 5, both ends of campaign & 0.5 \\
Processing & Registration, splits, datasheet, anonymisation & 3 \\
\midrule
\textbf{Total} & & \textbf{13.5} \\
\bottomrule
\end{tabular}
\end{center}

Against ten weeks to the CVPR deadline of 16 November 2026. The pilot is the only serial
dependency, so run it first.

\section{Quality gates}

Do not leave a floor until each of these passes. Re-surveying is far more expensive than
checking.

\begin{enumerate}
\item \textbf{Scale} --- x and y metres-per-pixel agree within 2\%, or the plan is
explicitly marked anisotropic.
\item \textbf{Grid closure} --- walking the grid perimeter returns to the origin within
0.3\,m.
\item \textbf{Coverage} --- every marked point has $\ge$7 device records, and the number
of distinct cells heard per point is within expectation for the floor.
\item \textbf{Non-detections present} --- spot-check that absent cells are recorded as
absent rather than omitted.
\item \textbf{Clock} --- the shared start event is visible in all seven logs and residual
offset is under 100\,ms after correction.
\item \textbf{Transmitter coverage} --- at least 15 physical units located per floor, each
mapped to its broadcast identifiers.
\item \textbf{Trajectory anchoring} --- every walk contains $\ge$4 marker presses and
starts and ends on a marked point.
\end{enumerate}

\section{Minimum viable, if time collapses}

Collect strictly in this order and stop wherever you must. Each prefix stands alone.

\begin{enumerate}
\item \textbf{CMU-Q f2 dense grid plus its transmitter survey.} This alone rescues the
real-data section and makes the geometry claim testable for the first time. 3 days.
\item \textbf{EC Parking f1, same two steps.} Adds the negative control, which is the
strongest single piece of evidence in the design. 2.5 days.
\item \textbf{CMU-Q f1.} Adds the cross-floor axis. 3 days.
\item Trajectories, inertial streams and device calibration across all three.
\end{enumerate}

Items 1 and 2 together, 5.5 days, deliver a defensible dataset contribution with a
treatment and a control. Everything after that widens it.

\end{document}
